%! Inverses of Exponential Functions — Unit 2, Lesson 2.10 — Group Activity (Answer Key)
\documentclass[10pt]{article}
\usepackage{precalculus-article}
\usepackage{precalculus-key}   % pulls in -boxes; do NOT also load -boxes
\newcommand{\TallMath}[1]{$\displaystyle #1\rule[-1.4em]{0pt}{3.2em}$}

\begin{document}

\pageheader{Unit 2, Lesson 2.10}{Group Activity --- Answer Key}
\namepartnerperiod

\vspace{0.1in}

\begin{scenariobox}[A Biologist's Inverse Problem]{sky}
A scientist models bacterial growth using $P(t) = 5^t$, where $P$ is the population count
and $t$ is time in hours. She needs the \emph{inverse} of this function to determine the time
elapsed given a current population count. Your group will construct and verify this inverse.
\end{scenariobox}

\vspace{0.1in}

\begin{tcolorbox}[colback=white, colframe=black!40,
    title=\textbf{Tier R --- Remediate},
    fonttitle=\bfseries, arc=2mm, left=3mm, right=3mm, top=2mm, bottom=2mm]

\textbf{Work with $g(x) = 4^x$.}

\begin{enumerate}[leftmargin=1.8em, itemsep=10pt]

\item Complete the table for $g(x) = 4^x$:

\vspace{4pt}
\begin{center}
\renewcommand{\arraystretch}{1.8}
\begin{tabular}{|c|c|}
\hline
\rowcolor{sky}
$x$ & $g(x) = 4^x$ \\
\hline
$0$ & \ans{1} \\
\hline
$1$ & \ans{4} \\
\hline
$2$ & \ans{16} \\
\hline
$3$ & \ans{64} \\
\hline
\end{tabular}
\end{center}

\item Swap the columns to make a table for the \emph{inverse} function $f(x) = \log_4 x$:

\vspace{4pt}
\begin{center}
\renewcommand{\arraystretch}{1.8}
\begin{tabular}{|c|c|}
\hline
\rowcolor{sky}
$x$ & $f(x) = \log_4 x$ \\
\hline
\ans{1}  & $0$ \\
\hline
\ans{4}  & $1$ \\
\hline
\ans{16} & $2$ \\
\hline
\ans{64} & $3$ \\
\hline
\end{tabular}
\end{center}

\item Using your table, evaluate: \quad $f(16) = \log_4 16 = $ \ans{2} \qquad $f(1) = \log_4 1 = $ \ans{0}

\vspace{6pt}

\item Verify the inverse: compute $g(f(16)) = 4^{\,f(16)} = 4^2 = $ \ans{16}. Does this equal 16? \ans{Yes}

\begin{teachernote}
Common error: students may write $\log_4 16 = 4$ (confusing the base with the answer).
Remind them: $\log_4 16$ asks ``to what power must we raise 4 to get 16?'' Answer: $4^2 = 16$,
so $\log_4 16 = 2$.
\end{teachernote}

\end{enumerate}
\end{tcolorbox}

\vspace{0.1in}

\begin{tcolorbox}[colback=white, colframe=black!40,
    title=\textbf{Tier A --- Approaching Proficiency},
    fonttitle=\bfseries, arc=2mm, left=3mm, right=3mm, top=2mm, bottom=2mm]

\textbf{Return to the scientist's model: $P(t) = 5^t$.}

\begin{enumerate}[leftmargin=1.8em, itemsep=10pt]

\item Write the inverse function $t(P)$ that gives hours elapsed given a population $P$.

\vspace{4pt}
$t(P) = $ \ans{$\log_5 P$}

\item Verify using composition. Show both directions:

\vspace{4pt}
$P(t(P)) = 5^{\log_5 P} = $ \ans{$P$} \quad $t(P(t)) = \log_5(5^t) = $ \ans{$t$}

\item How many hours until the population reaches 125?

\vspace{4pt}
$t(125) = $ \ans{$\log_5 125$} $= $ \ans{3} hours \quad (since $5^3 = 125$)

\item How many hours until the population reaches 625?

\vspace{4pt}
$t(625) = $ \ans{$\log_5 625$} $= $ \ans{4} hours \quad (since $5^4 = 625$)

\item Explain in words: how does the inverse function help the scientist?

\ansline{The model $P(t) = 5^t$ takes time as input and gives population as output. The
inverse $t(P) = \log_5 P$ reverses this: it takes population as input and gives elapsed time
as output, directly answering ``how long until we reach population $P$?''}

\end{enumerate}
\end{tcolorbox}

\vspace{0.1in}

\begin{tcolorbox}[colback=white, colframe=black!40,
    title=\textbf{Tier E --- Extension},
    fonttitle=\bfseries, arc=2mm, left=3mm, right=3mm, top=2mm, bottom=2mm]

\textbf{Comparing $f(x) = \log_2 x$ and $h(x) = 3\log_2 x$.}

\begin{enumerate}[leftmargin=1.8em, itemsep=10pt]

\item Build a table for both functions:

\vspace{4pt}
\begin{center}
\renewcommand{\arraystretch}{1.8}
\begin{tabular}{|c|c|c|}
\hline
\rowcolor{sky}
$x$ & $f(x) = \log_2 x$ & $h(x) = 3\log_2 x$ \\
\hline
$1$  & \ans{0} & \ans{0} \\
\hline
$2$  & \ans{1} & \ans{3} \\
\hline
$4$  & \ans{2} & \ans{6} \\
\hline
$8$  & \ans{3} & \ans{9} \\
\hline
$16$ & \ans{4} & \ans{12} \\
\hline
\end{tabular}
\end{center}

\item As $x$ doubles, $h(x)$ increases by \ans{3}. Is this additive growth?

\ansline{Yes. As $x$ increases multiplicatively (doubles), $h(x)$ increases additively (by 3
each time). This is the defining additive property of logarithmic functions.}

\item For what input(s) do $f(x) = h(x)$?

\ansline{Set $\log_2 x = 3\log_2 x$, so $0 = 2\log_2 x$, giving $\log_2 x = 0$, so $x = 1$.
They are equal only at $x = 1$ (both equal 0).}

\item Describe how the vertical dilation changes the outputs.

\ansline{Each output of $h(x)$ is 3 times the corresponding output of $f(x)$; the outputs
grow additively (by 3 per doubling) instead of 1 per doubling. $h(x)$ is still a logarithmic
function because it fits the general form $a\log_b x$ with $a = 3$, $b = 2$.}

\end{enumerate}
\end{tcolorbox}

\end{document}
